% Use \gls{key}, \Gls{key}, or \glspl{key} in the manuscript.
% Used entries appear automatically and terms link to their definitions.
\newglossaryentry{synthetic-data}{name={synthetic data},
  description={Artificially generated data designed to represent selected properties of observed or target data}}
% Abbreviations are displayed as short forms; the full name is in the glossary.
\newglossaryentry{machine-learning}{name={ML},
  description={Machine learning: methods that learn patterns from data to perform tasks such as prediction or classification}}
\newglossaryentry{gan}{name={GAN},plural={GANs},
  description={Generative adversarial network: a generative modelling approach in which a generator and a discriminator are trained with competing objectives}}
\newglossaryentry{ctgan}{name={CTGAN},
  description={Conditional tabular generative adversarial network: a GAN-based approach to generating synthetic tabular data using conditional generation}}
\newglossaryentry{fidelity}{name={fidelity},
  description={The extent to which synthetic data reproduce relevant properties of the target data}}
\newglossaryentry{utility}{name={utility},
  description={The usefulness of data for a specified downstream task and evaluation setting}}
